\chapter{Problemas de repaso}
\label{cap:problemas-repaso}

Los siguientes problemas integran definiciones, cálculos y referencias cruzadas. Las soluciones breves seleccionadas aparecen en el \cref{ap:soluciones}. En los demás, escribe el razonamiento completo antes de consultar otra fuente.

\begin{problem}
Usa \cref{eq:distancia-vertical} para calcular la distancia entre \(2i\) y \(5i\). ¿Qué isometría del semiplano puede trasladar ambos puntos verticalmente sin cambiar la respuesta?
\label{prob:verticales}
\end{problem}

\begin{problem}
En el disco de Poincaré, calcula la distancia entre el origen y \(z=1/3\), y verifica que el resultado coincide con \(\log 2\).
\label{prob:radial}
\end{problem}

\begin{problem}
Un triángulo tiene ángulos \(\pi/3\), \(\pi/4\) y \(\pi/6\). Calcula su área con \cref{eq:area-triangulo} y explica qué indica el signo del defecto angular.
\label{prob:area-triangulo}
\end{problem}

\begin{problem}
Decide si existe una teselación regular \(\{5,4\}\). Justifica la respuesta usando \cref{eq:condicion-teselacion}.
\label{prob:teselacion}
\end{problem}

\begin{problem}
Una superficie cerrada de género \(3\) tiene curvatura \(-1\). Calcula su área usando \cref{eq:area-superficie-cerrada}.
\label{prob:superficie}
\end{problem}

\begin{problem}
Elige uno de los resultados de los \cref{cap:disco-poincare,cap:semiplano-superior,cap:triangulos-hiperbolicos}. Reescribe su demostración o cálculo indicando cada definición usada y una fuente bibliográfica pertinente.
\label{prob:lectura}
\end{problem}

\section{Lista de comprobación}

Al terminar, comprueba que puedes explicar la diferencia entre coordenadas euclidianas e hiperbólicas, identificar las geodésicas de ambos modelos, calcular un defecto angular y usar una isometría para simplificar una configuración. El \cref{ap:guia-uso} muestra cómo seguir ampliando estas notas en LaTeX.
